\chapter{What finite reset experiments cannot exclude}
\label{rf:limits}

An observable that excludes a three-state classical model does not exclude
classical models with unrestricted hidden resources. This distinction survives
precise timestamps, adaptive scheduling, and arbitrarily accurate finite
records. This chapter constructs the alternative models explicitly. It also
shows that a transient coherent-looking signal need not identify nonzero
\emph{stationary} dissipation: a classical model satisfying detailed balance
can approximate the same records by starting far from equilibrium.

The approximation machinery is established. In particular, the Poisson
operator used below is the Sz\'asz--Mirakjan operator, introduced in the
classical approximation literature\cite{rfSzasz1950}. The role of the result
here is to locate the operational boundary of inference and count the hidden
resources needed by a sufficient countermodel. It is not a claim of a new
universal approximation principle.

\section{The access contract}

Let $a\in\{1,\ldots,m\}$ label finitely many preparations, and let
$p_a(t)\in[0,1]$ be the probability of a fixed binary detector outcome after
waiting for time $t$. Each record begins with a fresh preparation. Its arm and
waiting time may depend on earlier records, but are chosen independently of
the subsequent hidden noise in that record. The apparatus retains the arm,
waiting time, and outcome. The initial theorem concerns $t\in[0,T]$.

The classical comparison class has one common homogeneous finite Markov
generator, one fixed detector, and a distribution for each preparation. All
hidden clock and preparation-label states count toward its dimension. No
intermediate measurement or control during the wait is included. The theorem
does not reconstruct joint laws of repeated nondestructive measurements on a
single persistent system, nor does it match independently measured energy or
preparation-resource observables.

\section{A finite homogeneous classical realization}

\begin{theorem}[Explicit finite-window approximation]
\label{rf:poisson-realization}
Suppose all $p_a:[0,T]\to[0,1]$ have Lipschitz constant at most $L>0$.
For every $r>0$, there is a common homogeneous finite Markov generator with
deterministic binary readout, at most
\[
 n=2m(\lceil rT\rceil+1)
\]
states, and peak exit rate at most $r$, whose output probabilities satisfy
\begin{equation}
 \sup_{a,\,0\le t\le T}|q_a(t)-p_a(t)|
 \le L\sqrt{T/r}.
 \label{rf:lipschitz-realization}
\end{equation}
Thus $r=L^2T/\varepsilon^2$ is sufficient for error $\varepsilon$.
\end{theorem}

\begin{proof}
Set $K=\lceil rT\rceil$ and $f_a(k)=p_a(\min(k/r,T))$. A pure-birth clock
$k=0,\ldots,K$ starts at zero, jumps forward at rate $r$, and stays at $K$
thereafter. Introduce states $(a,k,y)$ with $y\in\{0,1\}$ and detector value
$y$. On moving to phase $k+1$, choose the new bit according to
$\operatorname{Bernoulli}(f_a(k+1))$. Initially choose the bit according to
$\operatorname{Bernoulli}(f_a(0))$. This is a homogeneous Markov model; the
label $a$ selects an initial component of its common generator.

If $N_{rt}$ is Poisson with mean $rt$, its detector probability is
\[
 q_a(t)=\mathbb E\,p_a\bigl(\min(N_{rt}/r,T)\bigr).
\]
Projection onto $[0,T]$ cannot increase distance from $t$. The Lipschitz
bound and Cauchy--Schwarz therefore give
\[
 |q_a(t)-p_a(t)|\le L\mathbb E|N_{rt}/r-t|
 \le L\sqrt{\operatorname{Var}(N_{rt}/r)}
 =L\sqrt{t/r}.
\]
The claimed state and rate bounds follow directly from the construction.
\end{proof}

Zero-probability bit states may be deleted. Uniform error $\eta$ in tabulating
the phase probabilities adds at most $\eta$ to the output bound. Preparations
with $p_a(0)=0$ or $1$ remain supported entirely on the appropriate detector
region: the construction does not rely on disguising an error at time zero.

\begin{proposition}[Smooth approximation and finite truncation]
\label{rf:smooth-realization}
Suppose $p_a:[0,\infty)\to[0,1]$ is twice differentiable and
$|p_a''(t)|\le M$ for every $a,t$. A clock frozen at phase $K$, tabulated using
$f_a(k)=p_a(k/r)$, satisfies
\begin{equation}
 \sup_{a,\,0\le t\le T}|q_a(t)-p_a(t)|
 \le \frac{MT}{2r}+\Pr\{\operatorname{Poisson}(rT)\ge K\}.
 \label{rf:smooth-clock-bound}
\end{equation}
\end{proposition}

\begin{proof}
For the untruncated clock, Taylor's linear term averages to zero because
$\mathbb E(N_{rt}/r)=t$. The second-order remainder is at most
$M\mathbb E(N_{rt}/r-t)^2/2=Mt/(2r)$. Freezing changes a bounded detector
mean by at most the probability of reaching the truncation region. Poisson
tail monotonicity in $t$ gives the stated bound.
\end{proof}

The global derivative assumption matters: the untruncated clock can query
times beyond $T$. For rational $\mu=rT$, an exact finite tail certificate is
\begin{equation}
 \Pr\{N\ge K\}\le
 \min_{0\le s\le K}\frac{\mu^s}{(K)_s},
 \qquad N\sim\operatorname{Poisson}(\mu),
 \label{rf:factorial-tail}
\end{equation}
where $(K)_s=K(K-1)\cdots(K-s+1)$ and the $s=0$ term is one.
It follows from the factorial moment $\mathbb E(N)_s=\mu^s$.
For very large state bounds the implementation also uses the simpler Cantelli
certificate
\[
 \Pr\{N\ge K\}\le \frac{\mu}{\mu+(K-\mu)^2},\qquad K>\mu.
\]
These are sufficient bounds, without an optimality claim.

For a finite-dimensional Lindblad generator $\mathcal L$ and fixed binary
effect $E$, all preparations satisfy usable global bounds
\[
 L=\|\mathcal L^*(E)\|_\infty,\qquad
 M=\|(\mathcal L^*)^2(E)\|_\infty,
\]
as chosen Lipschitz and curvature constants. Indeed the derivatives are expectations of these Hermitian
operators in the evolved density matrix. Thus every finite set of
preparations of a finite-dimensional Lindblad semigroup fits the construction.
The constructed classical clock rates need not equal the Lindblad damping
rates and are not inferred physical parameters.

\section{A detailed-balance countermodel with a preparation cost}

\begin{theorem}[Reversible realization]
\label{rf:reversible-realization}
Replace the clock's forward rate $r$ by forward and backward rates $r,d>0$,
with reflecting endpoints. Keep the destination-phase bit sampling used
above. The resulting classical model satisfies detailed balance and has zero
stationary entropy production. Its output probabilities differ from those of
the one-way clock by at most
\begin{equation}
 1-e^{-dT}\le dT
 \label{rf:backward-coupling}
\end{equation}
uniformly on $[0,T]$. Let $R=r/d$ and $Z=\sum_{k=0}^{K}R^k$. Relative to the
stationary law with equally weighted preparation components, the prescribed
initial state has relative entropy
\begin{equation}
 D(\nu_a\Vert\pi)=\log(mZ).
 \label{rf:preparation-cost}
\end{equation}
\end{theorem}

\begin{proof}
The phase chain has stationary weights $R^k/Z$. Write
$b_a(k,y)=f_a(k)^y[1-f_a(k)]^{1-y}$, with its usual Bernoulli interpretation
at endpoints. The lifted stationary law is
\[
 \pi(a,k,y)=\frac{R^k}{mZ}b_a(k,y).
\]
For neighboring phases $k,\ell$, the jump rate is $q_{k\ell}b_a(\ell,z)$.
The two stationary edge flows agree because
\[
 R^k b_a(k,y)q_{k\ell}b_a(\ell,z)
 =R^\ell b_a(\ell,z)q_{\ell k}b_a(k,y).
\]
Delete zero-weight bit states. Every retained state has positive stationary
weight within its component. Couple the forward jumps and sampled bits with
the one-way process until a backward jump occurs. The backward-jump counting
process is dominated by a rate-$d$ Poisson process, which proves
\eqref{rf:backward-coupling}. Finally,
$\nu_a(a,0,y)=b_a(0,y)$, so its likelihood ratio relative to $\pi$ equals
$mZ$ on its support. This proves \eqref{rf:preparation-cost}.
\end{proof}

If a single irreducible generator is required, add
$c(\mathbf 1\pi-I)$ for any $c>0$ on the retained states. This introduces
rare resets to $\pi$, connects all components, and preserves detailed balance.
A reset-event coupling adds at most $cT$ to the approximation error. Thus
arbitrarily accurate reversible approximations do not depend on reducibility
or absorbing states.

Zero \emph{stationary} entropy production does not imply zero transient
entropy production, zero reset cost, or free hardware. For an isothermal
detailed-balance realization at temperature $\Theta$,
$k_{\mathrm B}\Theta\log(mZ)$ is the nonequilibrium free-energy difference
from the specified equilibrium reference. It provides the corresponding
preparation-work lower bound from equilibrium when suitable reversible
controls are available. If the component label is conditioned on, the
within-component divergence is $\log Z$. The result assigns no universal
preparation charge to an arbitrary initial state supplied without a preparation
model. It does expose where a reversible classical imitation stores its
resource: highly nonequilibrium initialization of a large hidden clock.

\section{Retained timestamps, adaptive experiments, and a testing limit}

\begin{theorem}[Transcript approximation]
\label{rf:transcript-approximation}
Suppose the one-record output probabilities of the target and classical model
differ by at most $\varepsilon$ for every allowed preparation and waiting time.
For any fresh-reset protocol using at most $N$ records, including adaptive
choices and bounded stopping, their complete retained transcript laws obey
\begin{equation}
 \operatorname{TV}(P,Q)\le 1-(1-\varepsilon)^N\le N\varepsilon.
 \label{rf:transcript-tv}
\end{equation}
\end{theorem}

\begin{proof}
While the histories agree, share the random choices of next preparation and
waiting time. Conditional on those choices, couple the two Bernoulli outputs
maximally. The mismatch probability is at most $\varepsilon$ at every step.
Thus all records agree with probability at least $(1-\varepsilon)^N$.
The coupling inequality proves the result, including the retained timestamps.
\end{proof}

\begin{corollary}[No uniform quantum test against an unrestricted reset class]
\label{rf:closure-nogo}
Let $\mathcal C$ include every finite detailed-balance classical Markov model
with the access contract above and no bound on hidden dimension, rates, or
preparation resources. For a finite-record test $0\le\phi\le1$, suppose
\[
 \sup_{Q\in\mathcal C}\mathbb E_Q\phi\le\alpha.
\]
Then every smooth quantum prepare--wait--read law $P$ in the same contract
satisfies $\mathbb E_P\phi\le\alpha$.
\end{corollary}

\begin{proof}
For every $\varepsilon>0$, the preceding construction gives a classical
$Q_\varepsilon\in\mathcal C$ whose finite transcript law is within
$N\varepsilon$ of $P$. Expectations of a test valued in $[0,1]$ differ by at
most total variation. Therefore
$\mathbb E_P\phi\le\alpha+N\varepsilon$; let $\varepsilon$ tend to zero.
\end{proof}

The same statement applies to eventual finite rejection by a sequential test.
First restrict to rejection within $N$ records with every wait at most $T$;
then let $N,T$ increase. Every finite rejection is included, assuming each
chosen wait is finite and the target functions are regular on compact windows.
This argument permits the approximating classical model to grow. It does not
claim that a single fixed finite model matches infinitely many records.

The conclusion is a boundary on the null class, not a rejection of quantum
experiments. State-size, rate, preparation, energy, or intervention restrictions
can make a meaningful witness possible. Without such a restriction, the
unrestricted classical class is dense in the observed experiment. Increasing
the sample count alone cannot produce a test with controlled error and
nontrivial power against that entire class.

\section{A checked coherent-signal example}

For the damped qubit
\[
 p_0(t)=\frac{1+e^{-0.1t}\cos(2t)}2,\qquad p_1(t)=1-p_0(t),
 \qquad T=\frac94,
\]
a global curvature bound is $M=401/200$. The reproducible reversible
construction for error $1/20$ uses
\[
 r=\frac{10827}{80},\qquad d=\frac1{135},\qquad K=356.
\]
Its deterministic lift has 1,426 retained states, below the certified bound
1,428. The exact uniform error certificate is below $0.048250$; a numerical
101-point check finds maximum error approximately $0.009875$. The latter is
an illustration, while the former controls the full continuum of times.
The code checks 2,848 exact detailed-balance flow identities. Its preparation
relative entropy is approximately $3494.14$ nats, and its stationary entropy
production is zero.

For 6,000 records, a conservative Cantelli construction at per-record error
$1/120000$ gives a complete-transcript total-variation bound of $0.05$.
It uses $r=812025$, $d=1/810000$, and $K=2638068$, with a sufficient upper
bound of 10,552,276 deterministic hidden states. These are sufficient
resources, without a minimality or practical-hardware claim. They explain
precisely how a finite experiment can exclude a small classical model while
remaining compatible with a much larger equilibrium model supplied with a
nonequilibrium preparation. All of these checks concern synthetic models;
no laboratory measurements enter this example.

\section{A smaller exact realization and the discontinuity of exact inference}

The reversible approximation is a sufficient countermodel with a particular
thermodynamic property. It is not a minimum-state realization. For the damped
cosine, an exact and much smaller classical construction is available.
Consider a biased cycle with $n\ge3$ phases and $\theta=2\pi/n$. Choose forward
and backward rates $r,d$ by
\begin{equation}
 r+d=\frac{\gamma}{1-\cos\theta},\qquad
 r-d=\frac{\omega}{\sin\theta}.
 \label{rf:cyclic-rates}
\end{equation}
Its first Fourier mode has eigenvalue $-\gamma+i\omega$. From phase zero,
the fixed stochastic binary detector
\[
 f(k)=\frac{1+\cos(k\theta)}2
\]
therefore gives exactly
$[1+e^{-\gamma t}\cos(\omega t)]/2$ for every $t\ge0$. When $n$ is even,
preparation in phase $n/2$ gives the complementary trace.
Both rates are nonnegative precisely when
\begin{equation}
 n\arctan(\gamma/\omega)\ge\pi.
 \label{rf:cycle-size}
\end{equation}
These facts follow by applying the cycle generator to $e^{ik\theta}$.

For $\gamma=0.1$ and $\omega=2$, rational arctangent and Machin-formula bounds
certify that 62 fails \eqref{rf:cycle-size} and 63 satisfies it. Together with
the classical generator-sector lower bound, this gives an exact continuum
minimum of 63 states for the single trace when a fixed stochastic detector
is allowed. The 63-state rates are approximately $20.10533$ and $0.018529$.
A 64-state cycle realizes both complementary preparations; no minimality of
64 for that pair is asserted. If the detector must be deterministic within
the counted state space, destination-bit resampling supplies a lift with
125 states for the 63-cycle and 126 for the 64-cycle.

This model has stationary entropy production
$(r-d)\log(r/d)$ in units of $k_{\mathrm B}$ per unit time, and preparation
relative entropy $\log n$ from its uniform stationary phase
distribution. This law does not satisfy
detailed balance when $r\ne d$; no equilibrium thermodynamic interpretation
of that reference is implied here. The small exact cycle and the large
reversible approximation make different resource commitments.

Every finite detailed-balance generator has only real eigenvalues. It cannot
exactly produce the nonzero complex poles of the damped cosine. Exact agreement
on an open time interval would extend by analyticity. Nevertheless, the
reversible construction approximates this same signal arbitrarily closely on
every finite interval. Exact continuum exclusion and robust finite-precision
exclusion are therefore different inference problems. The latter requires
additional resource or access restrictions.
